\section{Reproducing the new polynomial constructions}
\label{sec:phase2-reproducibility}

The preferred new constructions are stored under
\nolinkurl{research/phase2/}. The complete paper is the standalone file
\nolinkurl{research/optimization/diophantine_optimization.tex}.
All variable counts include the distinguished input. The principal
specifications and their independent verification receipts are:
\begin{center}
\small
\begin{tabular}{r p{82mm}}
\toprule
Total variables & Package relative to \texttt{research/phase2/}\\
\midrule
43 & \nolinkurl{finite_fibre/polynomial-43-conic/}\\
45 & \nolinkurl{finite_fibre/polynomial-45-conic-semilocal/}\\
558 & \nolinkurl{finite_fibre/polynomial-558/}\\
1367 & \nolinkurl{quartic/quartic-best/}\\
1391 & \nolinkurl{quartic/quartic-small-coefficients/}\\
\bottomrule
\end{tabular}
\end{center}
The first three contain a
\nolinkurl{polynomial-specification.json}; the two quartics contain a
\nolinkurl{polynomial.json.gz}. Every package includes its construction
report and a \nolinkurl{verification.json} describing the exact scope
of its independent check. The complete SHA-256 manifest is generated
with the review archive. File names and hashes identify mathematical
objects; they are not additional mathematical hypotheses.

From the repository root, the single command
\begin{verbatim}
python3 -E -B research/phase2/verify_phase2.py
\end{verbatim}
runs eleven checks and records their outputs in
\nolinkurl{research/phase2/verification/all-checks.json} and its
\nolinkurl{logs/} directory. It checks the frozen source DAGs,
the scalar normalization constants, the conic parametrization and
duplication identities, the precision inequalities, the fixed-point
arithmetic, the three preferred finite-algebra specifications, and
both complete quartic tables. The checkers are separate programs from
the producers. Exact integer and rational arithmetic is used for the
new algebraic checks.

The finite-algebra interpreter also executes every determinant,
resultant, coefficient reversal, and outer finite case on eight small
examples, including nilpotents, excluded components, collisions,
rank drops, irrational residue fields, and zero quadratic roots. It
checks 177 factors and supplies a fully expanded degree-sixteen
example. This is a check of the construction on the specified examples.
The general theorem, and not extrapolation from those computations,
justifies the large instances.

To rebuild the quartics, use the producer with
\nolinkurl{--share-height --precision 52 --scalar-order 32}.
The optional \nolinkurl{--small-coefficients} selects the
1391-variable alternative. The package README files give the complete
commands for rebuilding each source and finite-algebra specification,
including all finite-chart and scalar-normalization options.
The original-source comparison checks remain independently runnable
using \nolinkurl{research/optimization/verify_all.py}.

These receipts establish the stated finite identities, inventories,
and bounds. They do not claim to verify the external prime-pattern
theorem, the source's restricted pole converse, or its full explicit
height theorem. Those arithmetic dependencies are stated in
Section~\ref{sec:scope} and in the source-audit appendix.
The huge indexed products and multiplication matrices are specified
exactly but are not materialized. Only the quartics have complete
expanded coefficient tables. No global minimum is claimed.

The assembly program checks balanced environments and braces, unique
labels, defined cross-references, and defined bibliography keys. A
working LaTeX format was not available in the preparation environment;
no successful typesetting or rendered-layout verification is claimed.
The standalone source can be compiled twice with a normal
\texttt{pdflatex} installation. All additions are prepared for review
on the local branch \nolinkurl{research/polynomial-optimization}.
No commit, push, or pull request has been made.
